\subsection{Additional identification details}\label{sec:appendix_identification}

First, we discuss the identification of unobserved heterogeneity. Unobserved types affect parameters of the perceived production functions, the utility from enrolling in college, and achievement. We discuss these sets of parameters separately.\par
\vspace{-10pt}
\paragraph{Type-dependent heterogeneity in beliefs.} Unobserved heterogeneity and measurement error on the survey answers used to elicit returns to effort generate variation across observationally identical students in perceived PSU scores, GPA, and returns to effort. We assume that the measurement error on the survey answers regarding hours of study under alternative hypothetical outcome scenarios, used to construct beliefs, is identically distributed to the measurement error on the reported actual hours of study. Therefore, variation in reported actual hours of study that is not explained by observed baseline characteristics identifies the variance of the measurement error. Having identified this parameter separately, we can use variation in beliefs between observationally identical students to pin down the unobserved heterogeneity in beliefs.\par 
\vspace{-10pt}
\paragraph{Type-dependent heterogeneity in the utility from enrolling in college.} Observationally identical students who face identical admission sets can make different enrollment decisions because of idiosyncratic preference shocks and because of permanent unobserved heterogeneity. To separately identify them we exploit the longitudinal aspect of our data. We observe student's preference-revealing choices at both the exam-taking decision stage and the enrollment stage. Unlike temporary preference shocks, permanent unobserved heterogeneity induces correlations in behavior over time, which allow us to pin down unobserved heterogeneity in the preference for college.\par 
\vspace{-10pt}
\paragraph{Type-dependent heterogeneity in achievement.} Observationally identical students can obtain different scores on the achievement test because of different type-dependent unobserved ability and different realizations of the measurement error. To separately identify them, first, we assume that the type is discrete and the measurement error is continuous. Therefore, the observed modes in the part of the achievement score not explained by observed characteristics are informative about type-specific ability. Second, we exploit the longitudinal aspect of our data. Students of different types obtain different achievement scores, exert different levels of effort, and make different educational choices. Unlike measurement error, permanent unobserved heterogeneity induces correlations between achievement, effort and later outcomes that are not explained by baseline characteristics and, therefore, are informative about unobserved heterogeneity.\par 

Second, we discuss how we mitigate potential endogeneity of the arguments of the subjective probability functions. For the subjective probability of a preferential admission, we use variation that comes from the experiment. The treatment makes this subjective probability salient: differences in choices across treatment groups are informative about the parameters of this subjective probability, because it governs pre-college behavior in the treatment group but not in the control group. For the subjective probability of a regular admission, we assume that there is a continuous characteristics (lagged achievement test score) that affects the expected entrance exam score but not the type distribution. Therefore, conditional on the variables that enter the type distribution (which include lagged GPA), variation in this lagged achievement score is exogenous. The intuition is that this variation captures idyosincratic, test-day shocks that are uncorrelated with a student's true ability or preferences. 


